% Strict LLNCS submission source generated from the locally compiled manuscript.
% Requires llncs.cls from the official template.
\documentclass[runningheads]{llncs}
\usepackage[T1]{fontenc}
\usepackage[utf8]{inputenc}
\usepackage{textcomp}
\usepackage{lmodern}
\usepackage{graphicx}
\usepackage{amsmath}
\usepackage{url}
\usepackage{array}
\usepackage{multirow}
\usepackage{makecell}
\usepackage{xcolor}
\usepackage{booktabs}


\newcommand{\safeincludegraphics}[2]{%
  \IfFileExists{#1}{\centering\includegraphics[width=0.88\textwidth]{#1}}{%
  \IfFileExists{../#1}{\centering\includegraphics[width=0.88\textwidth]{../#1}}{%
  \fbox{\parbox[c][0.22\textheight][c]{0.94\textwidth}{\centering\small #2\\[2mm]
  \footnotesize\url{#1}}}}}%
}
\providecommand{\doi}[1]{\url{https://doi.org/#1}}
\newcommand{\method}{GRAIN-CAN}
\newcommand{\framework}{ACE-CAN}
\newcommand{\attf}{F1$_{attack}$}
\newcommand{\wf}{F1$_{weighted}$}
\newcommand{\ratt}{R$_{attack}$}
\newcommand{\patt}{P$_{attack}$}

\setlength{\textfloatsep}{8pt plus 2pt minus 2pt}
\setlength{\floatsep}{8pt plus 2pt minus 2pt}
\setlength{\intextsep}{8pt plus 2pt minus 2pt}
\makeatletter
\setlength{\@fptop}{0pt}
\setlength{\@fpsep}{8pt plus 2pt minus 2pt}
\setlength{\@fpbot}{0pt plus 1fil}
\makeatother
\raggedbottom
\begin{document}
%
\title{Attack-Centric Evaluation and Feature-Preserving Detection for Cross-Vehicle Unknown-Attack CAN IDS}
\titlerunning{Attack-Centric CAN IDS Evaluation}
%
\author{Qi'ao Li}
\authorrunning{Qi'ao Li}
\institute{Anonymous Institution\\
\email{anonymous@example.com}}
%
\maketitle
%
\begin{abstract}
CAN intrusion detection is often evaluated on data where normal traffic overwhelms the attack class.  That is a poor fit for the job an IDS is supposed to do.  A detector has to find attacks at a tolerable false-alarm rate; recognizing benign traffic is not enough.  This problem is sharpest in cross-vehicle unknown-attack tests, where accuracy or weighted-F1 can be close to one even when no attack is detected.  We present \framework, an attack-centric evaluation protocol for CAN IDS.  It treats attacks as the positive class, audits normal-dominated scores, and reports attack precision, recall and F1, AUPR, Recall@FPR, false-alarm rates, trivial baselines, and event-level evidence.  We also introduce \method, a feature-preserving baseline that keeps same-ID timing, payload changes, and CAN-ID behavior at sample, short-window, and aggregate-window granularity.  On CT\&T test04, the cross-vehicle unknown-attack setting, \method{} reaches attack-F1 0.6678, AUPR 0.7845, and Recall@FPR=$10^{-3}$ of 0.8053.  With a validation-selected low-FPR threshold, the GRAIN window100 score stream reaches attack recall 0.8053 at actual test FPR 0.000681.  These results improve the corrected baseline, but they do not solve unknown-attack detection: default-threshold attack recall is 0.5189 and approximate event recall is 0.3650.  The evidence points to a simple reporting rule for CAN IDS papers: show attack-centric, low-FPR, and event-level results, together with class counts, threshold rules, and confusion matrices.
\keywords{Automotive security \and CAN bus \and intrusion detection \and class imbalance \and benchmark evaluation \and attack-centric metrics.}
\end{abstract}
%
\section{Introduction}
Modern vehicles rely on many electronic control units that exchange state over in-vehicle networks.  CAN is still widely used because it is simple, deterministic, and cheap.  It was not designed for hostile settings: it has no built-in message authentication, authorization, or confidentiality.  Once an attacker reaches the bus through an external interface or a compromised component, forged frames can affect vehicle functions or disrupt diagnostic traffic \cite{koscher2010experimental,checkoway2011comprehensive,miller2015remote}.  CAN IDS work therefore covers several kinds of evidence, including physical-layer fingerprints, timing, sequence models, and graph-based detectors \cite{cho2016fingerprinting,cho2017viden,lampe2023review,dupont2019network}.

The IDS objective is asymmetric.  A detector should raise alarms on attacks while keeping false alarms low.  It should not get credit for doing little more than naming the majority class.  CAN data make this easy to miss.  Normal frames are periodic and abundant, while attack frames can be rare, short, and tied to one vehicle or attack script.  A normal-dominated metric can therefore look excellent even when the detector misses every attack.  This is not a cosmetic metric issue.  A detector with near-perfect weighted-F1 and zero attack recall is not a deployable IDS.

We study this failure mode in cross-vehicle unknown-attack detection.  A detector trained on one vehicle and one attack family should not be assumed to work on a different vehicle and a different attack.  The label balance makes the problem worse.  In our CT\&T test04 sample-level audit, only 14,244 of 13,220,555 samples are attacks, a positive rate of 0.001077.  A classifier that predicts normal for every sample obtains 0.998923 accuracy and 0.998384 weighted-F1, with attack-F1 exactly zero.  An aggregate score can therefore look good while the detector finds no attacks.  CT\&T test04 is the case study, but the same failure can appear in any rare-attack benchmark.

We do not claim that a dataset is broken or that prior authors made a mistake.  The claim is narrower: in rare-attack CAN IDS, a metric table is hard to interpret unless it states the positive class, averaging rule, confusion matrix, threshold rule, and low-FPR behavior.  We use CT\&T test04 because it combines cross-vehicle shift, unknown attacks, and extreme imbalance.  We then ask a practical question: after the metrics are repaired, what kind of simple detector still works?

Our contributions are as follows.
\begin{itemize}
    \item \framework{} reports attack precision, recall and F1, AUPR, Recall@FPR, false-alarm rates, event-level evidence, trivial baselines, and ranking inversion for rare-attack CAN IDS.
    \item The metric forensics show how all-normal prediction can reach near-perfect accuracy and weighted-F1 while detecting no attacks.  Several Table-13-style rows also have reported recall equal to accuracy, which fits a weighted or accuracy-like interpretation.
    \item We build a corrected benchmark for four CT\&T settings and other data.  It covers trivial predictors, public-style ML, neural baselines, safe-feature protocols, and GRAIN rows.
    \item \method{} preserves same-ID timing, payload changes, payload statistics, and CAN-ID behavior.  It does not solve unknown attacks, but it is the best corrected baseline in our runs: on CT\&T test04 it reaches attack-F1 0.6678, AUPR 0.7845, AUROC 0.9024, and validation-threshold recall 0.8053 at actual FPR 0.000681.
    \item We add validation-threshold low-FPR evaluation, large-negative protocols, chunked-negative probes, and multi-seed feature-removal retraining.  The remaining gaps depend on external artifacts such as original code, confusion matrices, and official event boundaries.
\end{itemize}

\section{Background and First-Principles Motivation}
\subsection{CAN IDS Under Rare Attacks}
CAN messages expose only timestamps, arbitration IDs, data length codes, and up to eight payload bytes in classical CAN.  They usually do not expose signal names or ECU identities.  IDS methods therefore infer attacks from indirect evidence: timing regularity, ID frequency, transition structure, payload changes, or physical-layer traces.  CIDS and Viden use clock skew or voltage profiles to identify transmitters \cite{cho2016fingerprinting,cho2017viden}.  Data-driven methods learn temporal, geometric, graph, language-model, or statistical representations of frame streams \cite{tyree2018shape,verma2020road,alkhatib2022canbert,wang2023statgraph}.  ROAD, can-train-and-test, and can-sleuth show the same practical difficulty: public CAN data are limited, and results do not automatically transfer across vehicles or attack families \cite{verma2020road,lampe2024ctt,kidmose2025cansleuth}.

Rare-attack IDS evaluation should start with the attack class.  Missed attacks are the security failure, so attack-positive precision, recall and F1 must be visible.  Score metrics such as AUPR and Recall@FPR are needed because deployment uses false-alarm budgets.  Event-level analysis is also needed when attack intervals are available, since an operator cares whether an attack episode is detected and how quickly.  Accuracy and weighted-F1 can be useful context, but they cannot carry the detection claim.

\subsection{Why Weighted Metrics Can Mislead}
Let $N_0$ and $N_1$ be the numbers of normal and attack samples with $N_1 \ll N_0$.  A classifier that predicts all samples as normal has
\begin{equation}
    \mathrm{Accuracy}=\frac{N_0}{N_0+N_1}, \qquad F1_{attack}=0.
\end{equation}
When $N_1/(N_0+N_1)$ is below one percent, the accuracy is close to one even though the detector never detects an attack.  Weighted-F1 similarly aggregates class-wise F1 values according to support.  The normal class can dominate the final value, masking the zero attack-F1.  In single-label classification, weighted recall equals accuracy; therefore, a reported pattern in which recall and accuracy are identical is a warning sign that the metric may be support-weighted rather than attack-positive.

This observation is not specific to one algorithm.  It is a consequence of class imbalance and metric definition.  Any evaluation that does not report the attack-class confusion matrix can be misread.  The correct response is not to discard datasets, but to repair the evaluation protocol.

\subsection{Related Work and Positioning}
Automotive IDS research splits roughly into signal, protocol, and learning approaches.  Physical-layer methods fingerprint transmitters using clock skew, voltage characteristics, or other signal artifacts.  Cryptographic and authentication protocols address a different part of the problem by trying to prevent unauthenticated messages \cite{cho2016fingerprinting,cho2017viden,lu2019leap,lotto2024survey}.  Timing and statistical IDS methods use CAN periodicity, inter-arrival regularity, ID frequencies, and payload changes \cite{song2016intrusion,tyree2018shape}.

Neural IDS work uses recurrent models, transformers, and tokenized CAN frames \cite{taylor2016frequency,kang2016survival,alkhatib2022canbert}.  Graph learning and state-space models are also used for CAN traces \cite{wang2023statgraph,liu2026mids}.  Hardware IDS and UAV CAN detectors study related deployment settings \cite{seccan2025,khandelsurvey2025,majumder2024uavcan}.  We treat Mamba/MIDS-style and hardware IDS papers as related work rather than direct baselines because we do not have reproduced rows under the same CT\&T protocol.  Recent surveys and benchmark studies point to the same weak spot studied here: conclusions change with attack type, workload construction, protocol assumptions, and metric choice \cite{althunayyan2025survey,sharmin2024benchmarking,guerra2024road}.

Public datasets and benchmark studies, including ROAD, CT\&T, CrySyS, and can-sleuth, make comparison more feasible while exposing protocol mismatch and generalization difficulty \cite{verma2020road,lampe2024ctt,gazdag2023crysys,kidmose2025cansleuth}.  Imbalanced-evaluation work further motivates precision-recall analysis, MCC, and explicit positive-class reporting when attacks are rare \cite{davis2006pr,saito2015pr,chicco2020mcc}.

This paper is mostly about evaluation, but the detector still matters.  If a benchmark selects models by normal-dominated scores, a larger architecture may simply learn the wrong target.  If the metric suite is repaired but the representation discards same-ID timing and payload changes, the detector may still miss the attack.  \framework{} fixes the reporting target; \method{} gives the corrected benchmark a reproducible baseline.

\section{Attack-Centric Evaluation Framework}
\subsection{Problem Definition}
Let a CAN trace be an ordered sequence
\begin{equation}
    x_i=(t_i, id_i, dlc_i, b_{i,1},\ldots,b_{i,8}), \qquad y_i\in\{0,1\},
\end{equation}
where $y_i=1$ denotes an attack frame or an attack-derived window label.  A detector outputs either a binary decision $\hat{y}_i$ or a score $s_i$ where larger values indicate stronger attack evidence.  In cross-vehicle unknown-attack evaluation, the training distribution and test distribution differ in vehicle, attack type, or both.  The operational objective is not to maximize average label agreement; it is to select a threshold $\tau$ such that the detector recovers attack evidence while satisfying a false-positive budget:
\begin{equation}
    \mathrm{FPR}(\tau)=\frac{FP(\tau)}{FP(\tau)+TN(\tau)} \leq \alpha,
\end{equation}
for deployment budgets such as $\alpha=10^{-4}$ or $10^{-3}$.  This definition makes the attack class the positive class by construction.

\subsection{Metric Suite}
\framework{} reports results at three levels.
\paragraph{Sample/window classification.}  We report attack-positive precision, recall, and F1,
\begin{equation}
    P_{attack}=\frac{TP}{TP+FP}, \quad R_{attack}=\frac{TP}{TP+FN}, \quad
    F1_{attack}=\frac{2P_{attack}R_{attack}}{P_{attack}+R_{attack}},
\end{equation}
together with accuracy, macro-F1, weighted-F1, balanced accuracy, and MCC.  Accuracy and weighted-F1 are treated as diagnostic context, not as primary IDS metrics.
\paragraph{Score-based operation.}  For score-producing models, we report AUROC, AUPR, and Recall@FPR at $10^{-4}$, $5\times10^{-4}$, $10^{-3}$, $5\times10^{-3}$, and $10^{-2}$.  AUPR is included because precision-recall behavior is more informative than ROC behavior under severe imbalance \cite{davis2006pr,saito2015pr}.
\paragraph{Event-level evidence.}  When official attack intervals are unavailable, we construct transparent approximate event boundaries from contiguous positive segments within the same file/attack type and report the construction as a limitation.  We measure event recall, approximate detection delay, and false alarms per 100k samples or per hour when timestamps support it.

\subsection{Evaluation Protocol}
The corrected protocol makes the bookkeeping explicit.  Each dataset row reports positive and negative support before model scores.  Each binary metric states whether attacks are the positive class.  Thresholded rows state how the threshold was chosen; if only test score dumps are available, the row is marked as a best-test diagnostic rather than a validation-tuned deployment point.  The protocol also includes trivial baselines, because all-normal prediction is the quickest way to see whether a metric can be satisfied without attack detection.

\subsection{Metric Forensics and Ranking Inversion}
\framework{} also runs a small metric-forensics pass.  It evaluates all-normal and all-attack predictors, checks whether reported recall equals accuracy, compares reported F1 values with attack-positive, normal-positive, macro, weighted, and accuracy-like alternatives, and computes ranking inversion.  Ranking inversion asks a plain question: would the model selected by weighted-F1 or accuracy also be selected by attack-F1, AUPR, or Recall@FPR?  If the answer is no, the metric changes the engineering conclusion.

\begin{figure}[!htbp]
\safeincludegraphics{results/paper_revision_strict_review/figures/figure1_attack_centric_pipeline_fixed.pdf}{Attack-centric evaluation pipeline.  PDF conversion is required for camera-ready LLNCS builds.}
\caption{Attack-centric evaluation and feature-preserving detection pipeline.}
\label{fig:pipeline}
\end{figure}

\section{Feature-Preserving Detection Baseline}
\subsection{Why Granularity Matters}
A CAN attack may affect only a small number of frames.  A long window can hide local timing or payload changes, even when the downstream model is expressive.  Pure sample-level features have the opposite problem: they may miss short temporal patterns such as repeated insertions or period shifts.  GRAIN-CAN keeps the local causal signals and then aggregates them only after those signals have been computed.

\subsection{GRAIN-CAN Representation}
\method{} builds causal features from the current and past traffic only.  It does not use future frames, test labels, or global test-set statistics.  For each frame, it uses the most recent history of the same arbitration ID plus local stream context.  Window-level rows aggregate these causal frame features with summary statistics instead of replacing them with raw fixed-length token windows.  If an attack changes timing, payload continuity, or ID behavior, the representation should keep that change at the scale where it appears.

Its core feature groups are:
\begin{itemize}
    \item \textbf{Same-ID timing}: inter-arrival time and period deviation for the current arbitration ID.
    \item \textbf{Payload dynamics}: byte values, payload sum/standard deviation, and $L_1$ payload difference from the previous same-ID frame.
    \item \textbf{ID behavior}: CAN-ID frequency, entropy, and optional train-only transition rarity.
    \item \textbf{Granularity}: sample-level, short-window, and feature-preserving aggregate summaries.
\end{itemize}


\paragraph{Algorithm 1: causal GRAIN-CAN extraction.}
Given ordered CAN frames, a training split, a window size $k$, and a lightweight classifier:
(1) emit DLC/payload bytes, payload statistics, train-fit CAN-ID encodings, same-ID inter-arrival time, same-ID payload change, and causal ID-history summaries;
(2) fit all corpus-level encoders, frequency tables, scalers, and rarity summaries on training frames only;
(3) aggregate non-future window-$k$ groups by mean, standard deviation, minimum, maximum, and last value for timing, payload, and ID-behavior groups;
(4) assign a window attack label if at least one frame in the window is attack-labelled, never using labels as features; and
(5) train a score-producing classifier and report threshold-free AUPR/AUROC plus thresholded Recall@FPR.
The output is a feature-preserving representation plus a lightweight score-producing classifier; \method{} is not a new deep architecture.

We instantiate \method{} with lightweight tree classifiers.  The point is not to introduce another deep architecture.  The point is to test whether preserving CAN-specific evidence gives a better corrected baseline than fixed deep windows.  The method outputs scores for low-FPR analysis, supports feature and granularity audits, and keeps the model simple enough that the result is not hidden behind architecture complexity.

\subsection{Causality and Leakage Controls}
The representation uses only train-derived encoders or statistics when a feature requires corpus-level information.  Same-ID deltas are computed online from prior occurrences of the same ID.  Window labels are derived from attack presence in the window, but feature values do not use labels.  We separately audit shortcut-prone fields such as absolute timestamp protocols and distinguish safe CAN-derived features from risky metadata-derived features.  This does not eliminate every external-validity risk, but it makes the corrected benchmark a test of attack evidence rather than a test of file or timestamp leakage.

\section{Experimental Setup}
\subsection{Datasets and Settings}
We use CT\&T test01--test04 as the main cross-shift benchmark.  The settings separate known and unknown vehicles from known and unknown attack families.  Test04 is the hardest case in this paper because both the vehicle and the attack family shift.  ROAD, CrySyS subset/family data, HCRL, and Car-Hacking are used only as sanity checks for metric-trap behavior under different label balances and dataset conventions.

\begin{table}[!htbp]
\caption{Class imbalance and trivial all-normal behavior across settings.  The attack-F1 of all-normal prediction is always zero, but accuracy and weighted-F1 can be high under rare attacks.}
\label{tab:imbalance}
\centering
\resizebox{\textwidth}{!}{%
\begin{tabular}{lrrrrr}
\hline
Setting & Positives & Negatives & Pos. rate & All-normal acc. & All-normal \wf \\
\hline
ROAD & 9,352 & 272,840 & 0.03314 & 0.96686 & 0.95057 \\
CT\&T test01 & 1,128 & 162,419 & 0.00690 & 0.99310 & 0.98967 \\
CT\&T test02 & 2,137 & 168,863 & 0.01250 & 0.98750 & 0.98129 \\
CT\&T test03 & 704 & 192,168 & 0.00365 & 0.99635 & 0.99453 \\
CT\&T test04 & 638 & 238,087 & 0.00267 & 0.99733 & 0.99599 \\
CT\&T test04 sample & 14,244 & 13,206,311 & 0.00108 & 0.99892 & 0.99838 \\
CrySyS subset & 65,896 & 319,055 & 0.17118 & 0.82882 & 0.75124 \\
HCRL & 22,436 & 23,690 & 0.48641 & 0.51359 & 0.34855 \\
Car-Hacking & 23,314 & 152,262 & 0.13279 & 0.86721 & 0.80554 \\
\hline
\end{tabular}
}
\end{table}

\subsection{Baselines}
The baselines are grouped by purpose.  Trivial predictors expose metric traps.  Table-13-style classical models approximate public benchmark model families under corrected metrics.  Fixed-window neural baselines test whether generic sequence models recover cross-shift attacks.  Safe-feature and GRAIN representations test whether CAN timing and payload evidence should be retained across granularity.  The main rows are all-normal prediction, Table-13-style GradientBoosting/MLP, old window100 Transformer, safe-feature GradientBoosting, and GRAIN windows of 10, 20, and 100.  Earlier internal prototypes are excluded from the main comparison tables.

\subsection{Reproducibility and Claim Boundaries}
All headline claims are derived from repository-generated CSV/TEX/SVG artifacts under \texttt{results/attack\_centric\_final}, \texttt{results/final\_grain\_can}, and \texttt{results/final\_evidence\_completion}.  Event-level results are approximate if official event boundaries are unavailable.  The main GRAIN window100 low-FPR row now uses a validation-selected threshold; rows labelled best-test remain diagnostic score-separability upper bounds.  Feature-removal rows use completed capped-negative retraining over feature subsets on CT\&T test04.  They are stricter than proxy feature importance, but still do not replace exhaustive full-negative retraining.

\section{Results}
\subsection{RQ1: Do Aggregate Metrics Overstate Detection?}
Table~\ref{tab:imbalance} shows the failure directly.  In CT\&T test04 at sample level, an all-normal classifier obtains accuracy 0.99892 and weighted-F1 0.99838 while detecting no attacks.  At the aggregate setting used in the corrected benchmark, all-normal still obtains weighted-F1 0.99599 with attack-F1 0.  A near-perfect weighted-F1 score is therefore not evidence that the IDS detects attacks.

The Table-13 case study reinforces this point.  Nine rows in the public target table have reported accuracy equal to reported recall within $10^{-4}$.  Because weighted recall equals accuracy in single-label classification, this pattern is consistent with a weighted/accuracy-like metric hypothesis.  We do not infer author error; without confusion matrices and exact code, the safe conclusion is metric ambiguity.  But that ambiguity is enough to require attack-centric reporting.

\begin{figure}[!htbp]
\safeincludegraphics{results/paper_revision_strict_review/figures/figure2_table13_metric_forensics.pdf}{CT\&T Table 13 metric-forensics case study.  PDF conversion is required for camera-ready LLNCS builds.}
\caption{Metric forensics case study.  High aggregate scores under rare attacks can be consistent with normal-dominated metrics; they should not be interpreted as attack-positive F1 without confusion matrices.}
\label{fig:case}
\end{figure}

\subsection{RQ2: What Does the Corrected Benchmark Show?}
Table~\ref{tab:corrected} summarizes the best measured corrected rows over CT\&T test01--test04.  Test01 and test02 are largely recoverable under corrected metrics.  Test03 is harder.  Test04 is still the hardest row because both the vehicle and the attack family change.

\begin{table}[!htbp]
\caption{Best measured corrected rows on CT\&T settings.  Test04 remains substantially harder under attack-centric metrics.}
\label{tab:corrected}
\centering
\resizebox{\textwidth}{!}{%
\begin{tabular}{llrrrrr}
\hline
Setting & Best measured row & \attf & Prec. & Rec. & AUPR & AUROC \\
\hline
CT\&T test01 & GRAIN window 20 & 0.9685 & 0.9941 & 0.9442 & 0.9413 & 0.9723 \\
CT\&T test02 & GRAIN window 10 & 0.9948 & 0.9959 & 0.9937 & 0.9901 & 0.9967 \\
CT\&T test03 & GRAIN window 100 & 0.7967 & 0.9996 & 0.6623 & 0.9250 & 0.9598 \\
CT\&T test04 & GRAIN window 100 & 0.6678 & 0.9365 & 0.5189 & 0.7845 & 0.9024 \\
\hline
\end{tabular}
}
\end{table}

On test04, Table-13-style GradientBoosting reaches attack-F1 0.5477 and AUPR 0.3320.  The old window100 Transformer falls to attack-F1 0.0224 despite high recall because its precision is only 0.0113.  GRAIN window100 reaches attack-F1 0.6678 and AUPR 0.7845.  The tradeoff is still visible: the row is precise (\patt{} 0.9365), but it misses nearly half of the attack windows (\ratt{} 0.5189).  That supports feature preservation, but it does not support a claim that cross-vehicle unknown attacks are solved.

\begin{figure}[!htbp]
\safeincludegraphics{results/paper_revision_strict_review/figures/figure3_corrected_benchmark_heatmap.pdf}{Corrected CT\&T benchmark.  PDF conversion is required for camera-ready LLNCS builds.}
\caption{Corrected CT\&T benchmark heatmap.}
\label{fig:corrected}
\end{figure}

\subsection{RQ3: Does Metric Choice Change the Selected IDS?}
The ranking inversion analysis shows that metric choice can change the selected detector.  All-normal prediction is always last by attack-F1 but can rank well by weighted-F1 in highly imbalanced settings.  GRAIN-CAN ranks near the top under attack-centric metrics, especially in test02--test04.  For deployment, that matters: a normal-dominated metric can select a detector that a security engineer would reject.

\begin{figure}[!htbp]
\safeincludegraphics{results/paper_revision_strict_review/figures/figure4_ranking_inversion_scatter.pdf}{Ranking inversion between weighted-F1 and attack-F1.}
\caption{Ranking inversion between weighted-F1 and attack-F1.  All-normal prediction can look competitive under normal-dominated metrics but remains last under attack-F1.}
\label{fig:ranking}
\end{figure}

\subsection{RQ4: Why Does Feature Preservation Help?}
Table~\ref{tab:grain} compares test04 representations.  Sample-level detection reaches attack-F1 0.3884 and AUPR 0.2422.  Short windows improve the result, and feature-preserving window100 aggregation reaches attack-F1 0.6678 and AUPR 0.7845.  The old window100 deep Transformer obtains attack-F1 0.0224 and AUPR 0.1733.  The problem is not window length by itself.  The problem is whether the representation keeps the timing and payload evidence that survives the shift.  Table~\ref{tab:mechanism} reports feature-removal retraining evidence on CT\&T test04.  The study is capped-negative rather than full-negative, so it is mechanism evidence rather than a final ablation.

\begin{table}[!htbp]
\caption{Comparable GRAIN-CAN granularity results on CT\&T test04.}
\label{tab:grain}
\centering
\small
\begin{tabular*}{0.92\textwidth}{@{\extracolsep{\fill}}lrr@{}}
\hline
Representation & \attf & AUPR \\
\hline
Sample-level & 0.3884 & 0.2422 \\
GRAIN W10 & 0.6001 & 0.6185 \\
GRAIN W20 & 0.6416 & 0.7228 \\
GRAIN W100 aggregate & 0.6678 & 0.7845 \\
Old W100 Transformer & 0.0224 & 0.1733 \\
\hline
\end{tabular*}
\end{table}

\begin{figure}[!htbp]
\safeincludegraphics{results/paper_revision_strict_review/figures/figure5_grain_granularity.pdf}{GRAIN-CAN granularity comparison.}
\caption{Feature-preserving granularity on CT\&T test04.  Aggregate windows retain timing and payload evidence better than the old fixed-window deep representation in this setting.}
\label{fig:grain-granularity}
\end{figure}

\begin{table}[!htbp]
\caption{Feature-removal retraining evidence on CT\&T test04.}
\label{tab:mechanism}
\centering
\small
\begin{tabular*}{0.98\textwidth}{@{\extracolsep{\fill}}lrrrr@{}}
\hline
Variant & \attf & \ratt & AUPR & R@1e-3 \\
\hline
Full safe CAN & 0.4789 & 0.3540 & 0.2547 & 0.3550 \\
w/o $\Delta t_{sameID}$ & 0.0094 & 0.2486 & 0.0570 & 0.0668 \\
w/o payload delta & 0.4789 & 0.3540 & 0.2547 & 0.3550 \\
w/o payload stats & 0.4789 & 0.3540 & 0.2547 & 0.3550 \\
w/o CAN ID & 0.4797 & 0.3550 & 0.2649 & 0.3550 \\
Only timing & 0.5998 & 0.4909 & 0.4246 & 0.4909 \\
Only payload & 0.0016 & 0.0913 & 0.0135 & 0.0260 \\
Only ID & 0.0026 & 0.2555 & 0.0438 & 0.0550 \\
\hline
\end{tabular*}
\end{table}

The retraining rows point to same-ID timing as the dominant sample-level signal.  Removing $\Delta t_{sameID}$ collapses attack-F1, while the timing-only subset is strongest in this capped-negative study.  The three-seed summary shows the same pattern: timing-only reaches mean attack-F1 0.6013, while removing $\Delta t_{sameID}$ drops mean attack-F1 to 0.0085.  This matches the CAN setting: many attacks disturb periodicity, ID behavior, payload continuity, or short event structure.  A representation that averages these signals away will struggle even with an expressive classifier.

\subsection{RQ5: What Happens Under Low-FPR and Event-Level Metrics?}
Low-FPR metrics give a stricter view of deployment.  With validation-threshold evaluation, GRAIN window100 reaches AUPR 0.7868 on CT\&T test04.  A threshold selected on the validation split for an FPR budget of $5\times10^{-4}$ gives attack-F1 0.8639, attack recall 0.8053, and actual test FPR 0.000681.  The corresponding best-test Recall@FPR=$10^{-3}$ upper bound is 0.8053.  Safe-feature GradientBoosting improves over public-default sample protocols, but reaches only 0.3549 Recall@FPR=$10^{-3}$ and AUPR 0.2647.  The old window100 Transformer has high approximate event recall under its stored threshold, but weak AUPR and low-FPR recall, which points to poor score separation and false-positive pressure.  Best-test rows in Table~\ref{tab:lowfpr} are diagnostic score-separability evidence, not validation-tuned deployment thresholds.  Approximate event recall is built from labels, file continuity, and timestamps; it is not official event-boundary evidence.

\begin{table}[!htbp]
\caption{Low-FPR and approximate event evidence on CT\&T test04.  Validation rows use thresholds selected on the train-derived validation split; best-test rows are upper-bound score-separability diagnostics.}
\label{tab:lowfpr}
\centering
\resizebox{\textwidth}{!}{%
\begin{tabular}{llrrrr}
\hline
Model & Threshold & AUPR & Recall & Actual FPR & Event rec. \\
\hline
GRAIN window100 & val-FPR $5e{-}4$ & 0.7868 & 0.8053 & 0.000681 & 0.3650 \\
GRAIN window100 & best-test R@1e-3 & 0.7868 & 0.8053 & 0.000995 & 0.3650 \\
Safe-feature GB & best-test R@1e-3 & 0.2647 & 0.3549 & 0.001000 & 0.4758 \\
Old window100 Transformer & best-test R@1e-3 & 0.1732 & 0.1462 & 0.001000 & 1.0000 \\
Public-default HGB & best-test R@1e-3 & 0.0119 & 0.0109 & 0.001000 & 0.1637 \\
Predict all normal & default & 0.0027 & 0.0000 & 0.000000 & 0.0000 \\
\hline
\end{tabular}
}
\end{table}

\paragraph{Evidence closure and external sanity.}
Our audit separates local experiments from external dependencies.  Locally, we completed validation-threshold low-FPR rows, large-negative protocols for heavy classical baselines, chunked-negative probes, and multi-seed feature-removal retraining.  The remaining items, original confusion matrices/code, exact CT\&T-v1.5 alignment, official event boundaries, and exhaustive full-negative Protocol D, depend on external artifacts.  ROAD, CrySyS, HCRL, and Car-Hacking are used only as sanity checks for label balance and reporting risk.  They are not evidence of universal model dominance.  CT\&T test04 remains the main cross-vehicle unknown-attack case study.

\section{Discussion and Threats to Validity}
The negative result is part of the paper: cross-vehicle unknown-attack CAN IDS is not solved.  GRAIN window100 is the best corrected baseline in our runs, but its CT\&T test04 default-threshold recall and approximate event recall remain limited.  We also do not claim that CT\&T is wrong or that prior authors made an error.  Without original confusion matrices, exact code, and final metric settings, the defensible conclusion is metric ambiguity and evaluation risk.  Future CAN IDS papers should state attack-positive metrics, confusion matrices, threshold rules, trivial baselines, and event-level evidence where possible.

The benchmark correction is paired with a detector because metrics alone do not recover attack evidence.  A detector without corrected metrics can optimize the wrong target; a corrected benchmark without a detector leaves the practical question open.  The method result is intentionally plain: in these runs, preserving same-ID timing and payload dynamics mattered more than adding architecture complexity.  Future deep models can improve on \method{} if they encode the same invariants and use the same attack-centric protocol.

The limitations are concrete.  Some event-level rows use approximate boundaries from labels, files, attack types, and timestamp gaps rather than official attack intervals.  Some low-FPR rows are best-test diagnostic operating points because validation score dumps are unavailable.  ROAD, CrySyS, HCRL, and Car-Hacking are sanity checks for label balance and reporting risk, not proof of universal dominance.  Large-negative protocols are completed for the main classical model families, but exhaustive full-negative Protocol D remains separate.  The feature-removal study retrains feature subsets over multiple seeds, but it is still capped-negative.  The paper addresses evaluation and data-driven detection on CAN traces; it does not replace physical-layer authentication, cryptographic protocols, or system-level prevention.

\section{Conclusion}
Cross-vehicle unknown-attack CAN IDS should be evaluated from the attack class outward.  Accuracy and weighted-F1 can be dominated by normal traffic; in extreme rare-attack settings, an all-normal predictor can look nearly perfect while detecting no attacks.  \framework{} audits that failure mode, and \method{} provides a feature-preserving corrected baseline.  Some CT\&T settings are largely recoverable, but test04 remains difficult: the best measured GRAIN row reaches attack-F1 0.6678 and AUPR 0.7845, not an unknown-attack solution.  CAN IDS papers should report attack-positive, low-FPR, and event-level evidence with class distributions and confusion matrices before claiming robust detection.

\begin{credits}
\subsubsection{\ackname} This work was conducted as part of an academic study on automotive intrusion detection evaluation.  The authors thank the maintainers of public CAN IDS datasets and open-source tooling used for reproducible evaluation.

\subsubsection{\discintname} The authors have no competing interests to declare that are relevant to the content of this article.
\end{credits}

\begin{thebibliography}{30}
\bibitem{koscher2010experimental}
Koscher, K., Czeskis, A., Roesner, F., Patel, S., Kohno, T., Checkoway, S., McCoy, D., Kantor, B., Anderson, D., Shacham, H., Savage, S.: Experimental security analysis of a modern automobile. In: IEEE Symposium on Security and Privacy, pp. 447--462. IEEE (2010). \doi{10.1109/SP.2010.34}

\bibitem{checkoway2011comprehensive}
Checkoway, S., McCoy, D., Kantor, B., Anderson, D., Shacham, H., Savage, S., Koscher, K., Czeskis, A., Roesner, F., Kohno, T.: Comprehensive experimental analyses of automotive attack surfaces. In: USENIX Security Symposium. USENIX Association (2011)

\bibitem{cho2016fingerprinting}
Cho, K.T., Shin, K.G.: Fingerprinting electronic control units for vehicle intrusion detection. In: 25th USENIX Security Symposium, pp. 911--927. USENIX Association (2016)

\bibitem{cho2017viden}
Cho, K.T., Shin, K.G.: Viden: Attacker identification on in-vehicle networks. In: ACM SIGSAC Conference on Computer and Communications Security, pp. 1109--1123. ACM (2017). \doi{10.1145/3133956.3134001}

\bibitem{dupont2019network}
Dupont, G., den Hartog, J., Etalle, S., Lekidis, A.: Network intrusion detection systems for in-vehicle network. arXiv:1905.11587 (2019)

\bibitem{lampe2023review}
Lampe, B., Meng, W.: Intrusion detection in the automotive domain: a comprehensive review. IEEE Communications Surveys \& Tutorials \textbf{25}(4), 2356--2426 (2023). \doi{10.1109/COMST.2023.3309864}

\bibitem{tyree2018shape}
Tyree, Z., Bridges, R.A., Combs, F.L., Moore, M.R.: Exploiting the shape of CAN data for in-vehicle intrusion detection. In: IEEE Connected and Automated Vehicles Symposium, pp. 1--5. IEEE (2018)

\bibitem{verma2020road}
Verma, M.E., Bridges, R.A., Iannacone, M.D., Hollifield, S.C., Moriano, P., Hespeler, S.C., Kay, B., Combs, F.L.: A comprehensive guide to CAN IDS data and introduction of the ROAD dataset. arXiv:2012.14600 (2020)

\bibitem{lampe2024ctt}
Lampe, B., Meng, W.: can-train-and-test: a curated CAN dataset for automotive intrusion detection. Computers \& Security \textbf{140}, 103777 (2024). \doi{10.1016/j.cose.2024.103777}

\bibitem{kidmose2025cansleuth}
Kidmose, B., Kidmose, A., Meng, W.: can-sleuth: Sleuthing out the capabilities, limitations, and performance impacts of automotive intrusion detection datasets. International Journal of Information Security \textbf{24}, 193 (2025). \doi{10.1007/s10207-025-01038-8}

\bibitem{gazdag2023crysys}
Gazdag, A., Ferenc, R., Buttyan, L.: CrySyS dataset of CAN traffic logs containing fabrication and masquerade attacks. Scientific Data \textbf{10}, 903 (2023). \doi{10.1038/s41597-023-02716-9}

\bibitem{alkhatib2022canbert}
Alkhatib, N., Mushtaq, M., Ghauch, H., Danger, J.L.: CAN-BERT do it? Controller area network intrusion detection system based on BERT language model. arXiv:2210.09439 (2022)

\bibitem{wang2023statgraph}
Wang, K., Jiang, Q., Wang, B., Zhang, Y., Wu, Y.: Effective in-vehicle intrusion detection via multi-view statistical graph learning on CAN messages. arXiv:2311.07056 (2023)

\bibitem{guerra2024road}
Guerra, L., Xu, L., Bellavista, P., Chapuis, T., Duc, G., Mozharovskyi, P., Nguyen, V.T.: AI-driven intrusion detection systems on the ROAD dataset: a comparative analysis for automotive controller area network. arXiv:2408.17235 (2024)

\bibitem{taylor2016frequency}
Taylor, A., Leblanc, S., Japkowicz, N.: Anomaly detection in automobile control network data with long short-term memory networks. In: IEEE International Conference on Data Science and Advanced Analytics, pp. 130--139. IEEE (2016)

\bibitem{kang2016survival}
Kang, M.J., Kang, J.W.: Intrusion detection system using deep neural network for in-vehicle network security. PLOS ONE \textbf{11}(6), e0155781 (2016)

\bibitem{song2016intrusion}
Song, H.M., Kim, H.R., Kim, H.K.: Intrusion detection system based on the analysis of time intervals of CAN messages for in-vehicle network. In: International Conference on Information Networking, pp. 63--68. IEEE (2016)

\bibitem{miller2015remote}
Miller, C., Valasek, C.: Remote exploitation of an unaltered passenger vehicle. Black Hat USA (2015)

\bibitem{davis2006pr}
Davis, J., Goadrich, M.: The relationship between precision-recall and ROC curves. In: International Conference on Machine Learning, pp. 233--240. ACM (2006)

\bibitem{saito2015pr}
Saito, T., Rehmsmeier, M.: The precision-recall plot is more informative than the ROC plot when evaluating binary classifiers on imbalanced datasets. PLOS ONE \textbf{10}(3), e0118432 (2015). \doi{10.1371/journal.pone.0118432}

\bibitem{chicco2020mcc}
Chicco, D., Jurman, G.: The advantages of the Matthews correlation coefficient over F1 score and accuracy in binary classification evaluation. BMC Genomics \textbf{21}, 6 (2020). \doi{10.1186/s12864-019-6413-7}

\bibitem{liu2026mids}
Liu, Q., Song, R., Cui, L., Zhang, H., Sun, Y., Sun, L.: MIDS: Detecting stealthy masquerade and tampering attacks on CAN bus via bidirectional Mamba. arXiv:2606.18599 (2026)

\bibitem{lu2019leap}
Lu, Z., Wang, Q., Chen, X., Qu, G., Lyu, Y., Liu, Z.: LEAP: A lightweight encryption and authentication protocol for in-vehicle communications. arXiv:1909.10380 (2019)

\bibitem{lotto2024survey}
Lotto, A., Marchiori, F., Brighente, A., Conti, M.: A survey and comparative analysis of security properties of CAN authentication protocols. arXiv:2401.10736 (2024)

\bibitem{althunayyan2025survey}
Althunayyan, M., Javed, A., Rana, O.: A survey of learning-based intrusion detection systems for in-vehicle network. arXiv:2505.11551 (2025)

\bibitem{sharmin2024benchmarking}
Sharmin, S., Mansor, H., Abdul Kadir, A.F., Aziz, N.A.: Benchmarking frameworks and comparative studies of controller area network intrusion detection systems: a review. arXiv:2402.06904 (2024)

\bibitem{majumder2024uavcan}
Majumder, R., Comert, G., Werth, D., Gale, A., Chowdhury, M., Salek, M.S.: Graph-powered defense: Controller area network intrusion detection for unmanned aerial vehicles. arXiv:2412.02539 (2024)

\bibitem{seccan2025}
Khandelwal, S., Shanker, S.: SecCAN: An extended CAN controller with embedded intrusion detection. arXiv:2505.14924 (2025)

\bibitem{khandelsurvey2025}
Liu, Y., Xue, L., Wang, S., Luo, X., Zhao, K., Jing, P., Ma, X., Tang, Y., Zhou, H.: Vehicular intrusion detection system for controller area network: a comprehensive survey and evaluation. arXiv:2505.17274 (2025)

\end{thebibliography}
\end{document}
